% ================================================================
%  SESIÓN 12 — De enunciado a tabla
%  Almería en Cifras y Formas · Matemáticas 1.º ESO
%  IES Al-Ándalus · 2025-2026
% ================================================================
\documentclass[aspectratio=169, 10pt, spanish]{beamer}
\usepackage{almeria-beamer}

\renewcommand{\SessionNumber}{12}
\renewcommand{\SessionTitle}{De enunciado a tabla}
\renewcommand{\SessionName}{Sesión 12}
\renewcommand{\SessionObjective}{Organizar datos de un enunciado textual en una tabla ordenada y funcional, identificando las variables y sus valores}
\renewcommand{\BlockName}{Bloque 3: Datos y gráficas}

\begin{document}

% ---------------------------------------------------------------
%  PORTADA
% ---------------------------------------------------------------
\begin{frame}[plain]
  \titlepage
\end{frame}

% ---------------------------------------------------------------
%  FRAME 1 — ¿Por qué necesitamos tablas?
% ---------------------------------------------------------------
\begin{frame}{¿Por qué organizar datos en una tabla?}
  \begin{columns}[T]
    \begin{column}{0.52\textwidth}
      \begin{teoriabox}
        \AlmFontSmall
        Un \textbf{texto con datos} es difícil de analizar:
        \begin{itemize}
          \item Los números están dispersos entre palabras.
          \item No se aprecia el orden ni la magnitud.
          \item Es imposible hacer una gráfica directamente.
        \end{itemize}
        \vspace{4pt}
        Una \textbf{tabla} organiza la información de forma
        \textit{sistemática, ordenada y visual}.
      \end{teoriabox}
    \end{column}
    \begin{column}{0.44\textwidth}
      \begin{notabox}[Dato de Almería]
        \AlmFontSmall
        El Patronato de Turismo de Almería publica sus estadísticas
        en tablas mensuales.
        Sin esa organización, los datos serían ilegibles.
      \end{notabox}

      \vspace{0.3cm}
      \begin{center}
      \begin{tikzpicture}
        \node[font=\AlmFontLargeUp, color=AlmTerra] (texto) at (0,0)
          {\faFileAlt};
        \node[font=\AlmFontSmall, color=AlmTerra, below=2pt of texto]
          {Texto disperso};
        \draw[-Stealth, AlmAzulM, line width=2pt] (0.6,0) -- (1.6,0);
        \node[font=\AlmFontLargeUp, color=AlmVerde] (tabla) at (2.2,0)
          {\faTable};
        \node[font=\AlmFontSmall, color=AlmVerde, below=2pt of tabla]
          {Tabla ordenada};
      \end{tikzpicture}
      \end{center}
    \end{column}
  \end{columns}
\end{frame}

% ---------------------------------------------------------------
%  FRAME 2 — Estructura de una tabla
% ---------------------------------------------------------------
\begin{frame}{Estructura de una tabla de datos}
  \begin{center}
  \begin{tikzpicture}[font=\AlmFontSmall, every node/.style={align=center}]
    % Tabla visual con anotaciones
    \draw[AlmAzulM, line width=1.5pt] (0,0) rectangle (8,3.4);
    % Cabecera
    \fill[AlmAzul!85] (0,2.8) rectangle (8,3.4);
    \node[text=white, font=\bfseries\AlmFontSmall] at (2.4,3.1)
      {Encabezado col.~1};
    \draw[AlmAzulM, line width=1pt] (4.8,0) -- (4.8,3.4);
    \node[text=white, font=\bfseries\AlmFontSmall] at (6.4,3.1)
      {Encabezado col.~2};
    % Filas
    \foreach \y/\col in {2.2/AlmFondo, 1.6/white, 1.0/AlmFondo, 0.4/white}{
      \fill[\col, opacity=0.6] (0.1,\y-0.4) rectangle (7.9,\y+0.2);
    }
    \node at (2.4,2.1) {Fila 1 · celda A}; \node at (6.4,2.1) {Fila 1 · celda B};
    \node at (2.4,1.5) {Fila 2 · celda A}; \node at (6.4,1.5) {Fila 2 · celda B};
    \node at (2.4,0.9) {Fila 3 · celda A}; \node at (6.4,0.9) {Fila 3 · celda B};
    \node at (2.4,0.3) {\ldots};           \node at (6.4,0.3) {\ldots};
    % Anotaciones externas
    \draw[acota] (-0.7,2.8) -- (-0.7,3.4)
      node[midway, left=2pt, font=\AlmFontTiny\bfseries, color=AlmAzulM]
      {\rotatebox{90}{Encabezados}};
    \draw[acota] (-0.7,0) -- (-0.7,2.8)
      node[midway, left=2pt, font=\AlmFontTiny\bfseries, color=AlmGris]
      {\rotatebox{90}{Filas de datos}};
    \draw[<->, >=Stealth, AlmTerra, line width=0.8pt]
      (0,3.6) -- (4.7,3.6)
      node[midway, above, font=\AlmFontTiny\bfseries, color=AlmTerra]
      {Columna \(x\) independiente};
    \draw[<->, >=Stealth, AlmVerde, line width=0.8pt]
      (4.9,3.6) -- (8,3.6)
      node[midway, above, font=\AlmFontTiny\bfseries, color=AlmVerde]
      {Columna \(y\) dependiente};
  \end{tikzpicture}
  \end{center}

  \vspace{0.1cm}
  \begin{center}
  \AlmFontSmall
  \textbf{Regla:} columna izquierda = variable independiente \(x\);
  columna derecha = variable dependiente \(y\).
  \end{center}
\end{frame}

% ---------------------------------------------------------------
%  FRAME 3 — Pasos para construir una tabla
% ---------------------------------------------------------------
\begin{frame}{Pasos para construir una tabla desde un texto}
  \begin{columns}[T]
    \begin{column}{0.48\textwidth}
      \begin{pasobox}[1]
        \AlmFontSmall Lee con atención y \textbf{subraya todos los números}.
      \end{pasobox}

      \begin{pasobox}[2]
        \AlmFontSmall Identifica \textbf{qué representa cada número}
        (¿es una fecha? ¿una cantidad? ¿una medida?).
      \end{pasobox}

      \begin{pasobox}[3]
        \AlmFontSmall Decide cuál es la \textbf{variable independiente \(x\)}
        y cuál la \textbf{dependiente \(y\)}.
      \end{pasobox}
    \end{column}
    \begin{column}{0.48\textwidth}
      \begin{pasobox}[4]
        \AlmFontSmall Crea los \textbf{encabezados} con nombre claro y unidades
        entre paréntesis: \textit{Mes}, \textit{Visitantes (miles)}.
      \end{pasobox}

      \begin{pasobox}[5]
        \AlmFontSmall Rellena los valores en \textbf{orden creciente de \(x\)}
        (de menor a mayor).
      \end{pasobox}

      \begin{pasobox}[6]
        \AlmFontSmall \textbf{Verifica}: ¿están todos los datos?
        ¿son consistentes las unidades?
      \end{pasobox}
    \end{column}
  \end{columns}

  \vspace{0.2cm}
  \begin{notabox}[Clave]
    \AlmFontSmall
    Ordenar por \(x\) no es un capricho: facilita la lectura
    y es imprescindible para trazar una gráfica correcta después.
  \end{notabox}
\end{frame}

% ---------------------------------------------------------------
%  FRAME 4 — Ejemplo: visitantes a la Alcazaba
% ---------------------------------------------------------------
\begin{frame}{Ejemplo 1: visitantes a la Alcazaba}
  \begin{columns}[T]
    \begin{column}{0.50\textwidth}
      \begin{ejemplobox}
        \AlmFontSmall\itshape
        «En enero la Alcazaba recibió 8\,500 visitantes.
        En febrero, 9\,200. Marzo registró 12\,000 visitantes.
        En abril hubo 18\,500. En mayo la cifra llegó a 22\,000 personas.»
      \end{ejemplobox}

      \vspace{0.3cm}
      \begin{pasobox}[Análisis]
        \AlmFontSmall
        \begin{itemize}
          \item \(x\) = mes (cuantitativa discreta)
          \item \(y\) = visitantes (cuantitativa discreta)
          \item Ordenar: enero → mayo
        \end{itemize}
      \end{pasobox}
    \end{column}
    \begin{column}{0.46\textwidth}
      \begin{center}
      \renewcommand{\arraystretch}{1.35}
      \begin{tabular}{>{\bfseries\color{AlmAzulM}}c
                      >{\centering\arraybackslash}c
                      r}
        \toprule
        \rowcolor{AlmAzul!12}
        \textbf{Mes} & \textbf{N.º mes} & \textbf{Visitantes} \\
        \midrule
        Enero    & 1 & 8\,500  \\
        \rowcolor{AlmFondo}
        Febrero  & 2 & 9\,200  \\
        Marzo    & 3 & 12\,000 \\
        \rowcolor{AlmFondo}
        Abril    & 4 & 18\,500 \\
        Mayo     & 5 & 22\,000 \\
        \bottomrule
      \end{tabular}
      \end{center}
      \vspace{0.2cm}
      \begin{notabox}[Observa]
        \AlmFontSmall
        ¡Los datos estaban desordenados en el texto!
        La tabla los ordena por el mes.
      \end{notabox}
    \end{column}
  \end{columns}
\end{frame}

% ---------------------------------------------------------------
%  FRAME 5 — Ejemplo 2: tarifa de taxi en Almería
% ---------------------------------------------------------------
\begin{frame}{Ejemplo 2: tarifa de taxi en Almería}
  \begin{columns}[T]
    \begin{column}{0.50\textwidth}
      \begin{ejemplobox}
        \AlmFontSmall\itshape
        «La bajada de bandera es 2{,}40\,€.
        Cada kilómetro recorrido cuesta 1{,}20\,€ adicionales.
        Calcula el precio para 1, 2, 3, 4 y 5\,km.»
      \end{ejemplobox}

      \vspace{0.25cm}
      \begin{formulabox}
        $y = 2{,}40 + 1{,}20 \cdot x$
      \end{formulabox}

      \vspace{0.2cm}
      {\AlmFontSmall $x$ = km recorridos \quad $y$ = precio total (€)}
    \end{column}
    \begin{column}{0.46\textwidth}
      \begin{center}
      \renewcommand{\arraystretch}{1.35}
      \begin{tabular}{>{\bfseries\color{AlmAzulM}\centering\arraybackslash}p{1.8cm}
                      >{\centering\arraybackslash}p{2.4cm}}
        \toprule
        \rowcolor{AlmAzul!12}
        \textbf{Km (\(x\))} & \textbf{Precio €\,(\(y\))} \\
        \midrule
        1 & 2{,}40 + 1{,}20 = 3{,}60 \\
        \rowcolor{AlmFondo}
        2 & 2{,}40 + 2{,}40 = 4{,}80 \\
        3 & 2{,}40 + 3{,}60 = 6{,}00 \\
        \rowcolor{AlmFondo}
        4 & 2{,}40 + 4{,}80 = 7{,}20 \\
        5 & 2{,}40 + 6{,}00 = 8{,}40 \\
        \bottomrule
      \end{tabular}
      \end{center}
    \end{column}
  \end{columns}
\end{frame}

% ---------------------------------------------------------------
%  FRAME 6 — Tabla de frecuencias
% ---------------------------------------------------------------
\begin{frame}{Tabla de frecuencias: cuando muchos datos se repiten}
  \begin{columns}[T]
    \begin{column}{0.50\textwidth}
      \begin{teoriabox}
        \AlmFontSmall
        Cuando hay \textbf{muchas repeticiones} del mismo valor,
        usamos una \textbf{tabla de frecuencias}:
        \begin{itemize}
          \item \textbf{Valor}: cada valor posible de la variable.
          \item \textbf{Marcas de conteo}: palotes para contar
                (5 palotes = \texttt{IIII\hspace{-1pt}/}).
          \item \textbf{Frecuencia}: número total de repeticiones.
        \end{itemize}
      \end{teoriabox}
    \end{column}
    \begin{column}{0.46\textwidth}
      \begin{center}
      \AlmFontSmall
      \textit{Ejemplo: tipo de transporte en Almería (encuesta a 30 alumnos)}\\[4pt]
      \renewcommand{\arraystretch}{1.25}
      \begin{tabular}{lcc}
        \toprule
        \rowcolor{AlmAzul!12}
        \textbf{Transporte} & \textbf{Marcas} & \textbf{Frec.} \\
        \midrule
        A pie      & \texttt{IIIII IIIII II}  & 12 \\
        \rowcolor{AlmFondo}
        Autobús    & \texttt{IIIII IIII}       &  9 \\
        Coche      & \texttt{IIIII I}          &  6 \\
        \rowcolor{AlmFondo}
        Bicicleta  & \texttt{III}              &  3 \\
        \midrule
        \textbf{Total} & & \textbf{30} \\
        \bottomrule
      \end{tabular}
      \end{center}
    \end{column}
  \end{columns}
\end{frame}

% ---------------------------------------------------------------
%  FRAME 7 — Gráfico de barras a partir de la tabla (pgfplots)
% ---------------------------------------------------------------
\begin{frame}{De la tabla al gráfico: Alcazaba enero–mayo}
  \begin{center}
  \begin{tikzpicture}
    \begin{axis}[
      almeria bar,
      width=0.82\linewidth, height=0.56\linewidth,
      title={Visitantes mensuales a la Alcazaba},
      xlabel={Mes},
      ylabel={Visitantes},
      symbolic x coords={Enero,Febrero,Marzo,Abril,Mayo},
      xtick=data,
      ymin=0, ymax=26000,
      nodes near coords style={font=\AlmFontTiny\bfseries, color=AlmAzul},
      bar width=0.6cm,
    ]
      \addplot[fill=AlmAzulM, draw=AlmAzul] coordinates {
        (Enero,8500)
        (Febrero,9200)
        (Marzo,12000)
        (Abril,18500)
        (Mayo,22000)
      };
    \end{axis}
  \end{tikzpicture}
  \end{center}
  \vspace{-0.2cm}
  {\AlmFontSmall\color{AlmGris} La tabla se convierte directamente en este gráfico
  gracias al orden correcto de los datos.}
\end{frame}

% ---------------------------------------------------------------
%  FRAME 8 — Ejercicio RECORDAR
% ---------------------------------------------------------------
\begin{frame}{Ejercicio: elementos de una tabla}
  \begin{recordarbox}
    Responde sin mirar los apuntes:
    \begin{enumerate}
      \item ¿Cuáles son los \textbf{4 elementos} que debe tener
            toda tabla bien estructurada?
      \item ¿Qué información debe aparecer en los
            \textbf{encabezados de columna}?
      \item ¿Qué variable va en la columna \textbf{izquierda}?
            ¿Y en la derecha?
      \item ¿Por qué hay que ordenar los datos
            de menor a mayor valor de \(x\)?
    \end{enumerate}
  \end{recordarbox}

  \vspace{0.3cm}
  \begin{center}
  \AlmFontSmall
  \begin{tabular}{>{\bfseries\color{AlmAzulM}}l l}
    \toprule
    Elemento & Descripción \\
    \midrule
    Encabezados & Nombre de variable + unidad \\
    Filas & Un dato por fila \\
    Columna \(x\) & Variable independiente (izquierda) \\
    Columna \(y\) & Variable dependiente (derecha) \\
    \bottomrule
  \end{tabular}
  \end{center}
\end{frame}

% ---------------------------------------------------------------
%  FRAME 9 — Ejercicio COMPRENDER
% ---------------------------------------------------------------
\begin{frame}{Ejercicio: detectar errores en una tabla}
  \begin{comprenderbox}
    La siguiente tabla tiene \textbf{4 errores}. Identifícalos
    y reescribe la tabla de forma correcta.

    \vspace{0.25cm}
    \begin{center}
    \AlmFontSmall
    \begin{tabular}{cc}
      \toprule
      \rowcolor{AlmTerra!10}
      \textbf{Datos} & \textbf{Cosas} \\
      \midrule
      Julio  & 280000 \\
      Enero  & 60 \\
      Marzo  & 80000 \\
      Agosto & mucho \\
      Febrero& 65000 \\
      \bottomrule
    \end{tabular}
    \end{center}

    \vspace{0.15cm}
    \textit{Pista}: busca problemas en encabezados, unidades,
    orden y consistencia de datos.
  \end{comprenderbox}
\end{frame}

% ---------------------------------------------------------------
%  FRAME 10 — Ejercicio APLICAR
% ---------------------------------------------------------------
\begin{frame}{Ejercicio: el autobús de Almería}
  \begin{aplicarbox}
    Lee el enunciado y construye la tabla correspondiente:

    \vspace{0.2cm}
    \begin{center}
    \begin{tcolorbox}[colback=AlmAmbar!12, colframe=AlmAmbar!60!black,
                      width=0.88\textwidth, arc=4pt, boxrule=0.8pt]
      \AlmFontSmall\itshape
      «El autobús urbano de Almería tarda 5 minutos en hacer 2 paradas.
      Para 4 paradas necesita 10 minutos.
      Con 6 paradas son 15 minutos.
      Las 8 paradas se completan en 20 minutos,
      y las 12 paradas en 30 minutos.»
    \end{tcolorbox}
    \end{center}

    \vspace{0.2cm}
    \begin{enumerate}
      \item Construye la tabla completa con encabezados y unidades.
      \item ¿Cuál es la variable independiente? ¿Y la dependiente?
      \item ¿Ves un patrón? ¿Cuánto tarda para \textbf{10 paradas}?
    \end{enumerate}
  \end{aplicarbox}
\end{frame}

% ---------------------------------------------------------------
%  FRAME 11 — Error frecuente
% ---------------------------------------------------------------
\begin{frame}{¡Cuidado con este error!}
  \begin{errorbox}
    \textbf{Meter los datos en orden aleatorio}:

    \vspace{0.3cm}
    \begin{columns}[T]
      \begin{column}{0.44\textwidth}
        \incorrecto\; \textbf{Tabla desordenada}
        \vspace{4pt}
        \AlmFontSmall
        \begin{tabular}{cc}
          \toprule
          \rowcolor{AlmTerra!10}
          Mes & Visitantes \\
          \midrule
          Julio & 500\,000 \\
          Enero & 60\,000 \\
          Abril & 120\,000 \\
          Marzo & 80\,000 \\
          \bottomrule
        \end{tabular}
      \end{column}
      \begin{column}{0.44\textwidth}
        \correcto\; \textbf{Tabla ordenada}
        \vspace{4pt}
        \AlmFontSmall
        \begin{tabular}{cc}
          \toprule
          \rowcolor{AlmVerde!10}
          Mes & Visitantes \\
          \midrule
          Enero & 60\,000 \\
          Marzo & 80\,000 \\
          Abril & 120\,000 \\
          Julio & 500\,000 \\
          \bottomrule
        \end{tabular}
      \end{column}
    \end{columns}

    \vspace{0.25cm}
    \AlmFontSmall
    Una tabla desordenada \textbf{impide ver la tendencia} y genera gráficas
    con puntos conectados de forma caótica.
    Ordena siempre de menor a mayor valor de \(x\).
  \end{errorbox}
\end{frame}

% ---------------------------------------------------------------
%  FRAME 12 — Evidencia del día
% ---------------------------------------------------------------
\begin{frame}{Evidencia del día}
  \begin{evidenciabox}
    Construye en tu cuaderno \textbf{dos tablas completas}
    a partir de estos dos párrafos:

    \vspace{0.2cm}
    \begin{tcolorbox}[colback=AlmAmbar!10, colframe=AlmAmbar!60!black,
                      arc=3pt, boxrule=0.8pt, fonttitle=\bfseries\AlmFontSmall,
                      title={Párrafo A — Temperaturas en Almería}]
      \AlmFontSmall\itshape
      «En enero la temperatura media es 12\,°C.
      En abril sube a 19\,°C. En julio alcanza 30\,°C.
      Octubre registra 21\,°C y diciembre baja a 13\,°C.»
    \end{tcolorbox}

    \vspace{0.2cm}
    \begin{tcolorbox}[colback=AlmAzulC!10, colframe=AlmAzulC!70!black,
                      arc=3pt, boxrule=0.8pt, fonttitle=\bfseries\AlmFontSmall,
                      title={Párrafo B — Precio del taxi}]
      \AlmFontSmall\itshape
      «Un taxi en Almería cobra 1{,}80\,€/km con bajada de bandera de 3\,€.
      Calcula el precio total para 1, 3, 5, 7 y 10\,km.»
    \end{tcolorbox}

    \vspace{0.2cm}
    \AlmFontSmall
    Cada tabla debe tener: \textbf{encabezados con unidades},
    \textbf{orden correcto de \(x\)} y \textbf{5 filas de datos}.
  \end{evidenciabox}
\end{frame}

% ---------------------------------------------------------------
%  FRAME 13 — Resumen de sesión
% ---------------------------------------------------------------
\begin{frame}{Resumen: lo que hemos aprendido hoy}
  \begin{resumenbox}
    \begin{columns}[T]
      \begin{column}{0.50\textwidth}
      \vspace{0.15cm}
        \textbf{Estructura de una tabla:}
        \begin{itemize}
          \item Encabezados con nombre y unidades
          \item Columna izquierda = \(x\) independiente
          \item Columna derecha = \(y\) dependiente
          \item Filas ordenadas de menor a mayor \(x\)
        \end{itemize}
      \end{column}
      \begin{column}{0.46\textwidth}
        \textbf{6 pasos para construir la tabla:}
        \begin{enumerate}
          \AlmFontSmall
          \item Subraya todos los números
          \item Identifica qué representa cada uno
          \item Decide \(x\) e \(y\)
          \item Crea encabezados claros
          \item Rellena en orden
          \item Verifica consistencia
        \end{enumerate}
      \end{column}
    \end{columns}

    \vspace{0.2cm}
    \separador
    \vspace{0.15cm}
    \AlmFontSmall
    \textbf{Error a evitar:}
    datos en orden aleatorio dificultan la lectura y hacen
    imposible trazar una gráfica correcta.
  \end{resumenbox}

  \vspace{0.2cm}
  \begin{center}
    \bloomtag{Recordar}\quad
    \bloomtag{Comprender}\quad
    \bloomtag{Aplicar}
  \end{center}
\end{frame}

\end{document}


